\begin{longtable}[c]{|p{3.5cm}|p{12cm}|} 
    
  \caption{Data Frames and Remote Frames in SFF and EFF}
  \label{tab:twai-data-frames-and-remote-frames} 
  \\\hline
  \rowcolor{lightgray}
    \textbf{Data/Remote Frames} & \textbf{Description}
  \\\hline
  \endfirsthead
  
  \multicolumn{2}{c}{\bfseries \tablename\ \thetable{} -- cont'd from previous page}\\  \hline
  
  \rowcolor{lightgray}
    \textbf{Data/Remote Frames} & \textbf{Description}
  \\\hline
  \endhead
  
  \multicolumn{2}{r}{\textbf{Cont'd on next page}} \\
  \endfoot
  \endlastfoot
  
    SOF & The SOF (Start of Frame) is a single dominant bit used to synchronize nodes on the bus.
          \\\hline
    Base ID & The Base ID (ID.28 to ID.18) is the 11-bit identifier for SFF, or the first 11 bits of the 29-bit identifier for EFF.
          \\\hline
    \label{table_item:canframe_rtr}{RTR} & The RTR (Remote Transmission Request) bit indicates whether the message is a data frame (dominant) or a remote frame (recessive). This means that a remote frame will always lose arbitration to a data frame if they have the same ID.                                                    \\\hline
    \label{table_item:canframe_srr}{SRR} & The SRR (Substitute Remote Request) bit is transmitted in EFF to substitute for the RTR bit at the same position in SFF.                       
           \\\hline
    IDE & The IDE (Identifier Extension) bit indicates whether the message is SFF (dominant) or EFF (recessive). This means that a SFF frame will always win arbitration over an EFF frame if they have the same Base ID.
            \\\hline
    Extd ID & The Extended ID (ID.17 to ID.0) is the remaining 18 bits of the 29-bit identifier for EFF.
            \\\hline
    r1 & The r1 bit (reserved bit 1) is always dominant.
             \\\hline 
    r0 & The r0 bit (reserved bit 0) is always dominant.
             \\\hline
    DLC & The DLC (Data Length Code) is 4-bit long and should contain any value from 0 to 8. Data frames use the DLC to indicate the number of data bytes in the data frame. Remote frames used the DLC to indicate the number of data bytes to request from another node. 
             \\\hline
    Data Bytes & The data payload of data frames. The number of bytes should match the value of DLC. Data byte 0 is transmitted first, and each data byte is transmitted from the most significant bit first.
             \\\hline
    CRC Sequence & The CRC sequence is a 15-bit cyclic redundancy code.
             \\\hline
    CRC Delim & The CRC Delim (CRC Delimiter) is a single recessive bit that follows the CRC sequence.
              \\\hline
    ACK Slot & The ACK Slot (Acknowledgment Slot) is intended for receiver nodes to indicate that the data or remote frame was received without any issue. The transmitter node will send a recessive bit in the ACK Slot and receiver nodes should override the ACK Slot with a dominant bit if the frame was received without errors.
              \\\hline
    ACK Delim & The ACK Delim (Acknowledgment Delimiter) is a single recessive bit.
              \\\hline
    EOF & The EOF (End of Frame) marks the end of a data or remote frame, and consists of seven recessive bits.
              \\\hline
\end{longtable}
